\begin{tikzpicture}[
    node distance=0.8cm and 3cm,
    every node/.style={font=\small},
    block/.style={rectangle, draw, fill=blue!10, text width=5em, text centered, rounded corners, minimum height=2.5em},
    decision/.style={diamond, draw, fill=yellow!10, text width=4em, text centered, aspect=2, inner sep=1pt},
    action/.style={rectangle, draw, fill=green!10, text width=6em, text centered, rounded corners, minimum height=2.5em},
    suspend/.style={rectangle, draw, fill=orange!20, text width=6em, text centered, rounded corners, minimum height=2.5em},
    title/.style={font=\large\bfseries, text width=13em, align=center},
    note/.style={rectangle, draw, fill=gray!10, text width=6em, text centered, rounded corners, minimum height=2.5em},
    whiteNote/.style={rectangle, draw, fill=white, text width=6em, text centered, rounded corners, minimum height=2.5em},
    pinkNote/.style={rectangle, draw, fill=orange!10, text width=6em, text centered, rounded corners, minimum height=2.5em},
    bgBox/.style={rectangle, draw=gray!50, fill=orange!10, rounded corners, inner sep=0.5cm},
    container/.style={rectangle, draw, fill=orange!10, text width=8em, minimum height=6em, rounded corners},
    arrow/.style={->, >=stealth, thick},
    transfer/.style={->, >=stealth, thick, dashed},
    solidarrow/.style={->, >=stealth, thick},
    label/.style={font=\scriptsize}
]
\pgfdeclarelayer{background}
\pgfsetlayers{background,main}
% -------------------------
% Column Titles & Alignment Coordinates
% -------------------------
    \node[title] (t1) at (0, 0) {\texttt{startCoroutine-\\UninterceptedOrReturn}};
    \node[title] (t2) at (6, 0) {\texttt{CoroutineImpl::\\resumeWith}};
    \node[title] (t3) at (12, 0) {suspend\\function};
    \node[title] (t4) at (18, 0) {\texttt{suspendCoroutine-\\UninterceptedOrReturn}};

% -------------------------
% Column 1 Nodes
% -------------------------
    \node[whiteNote, text width=7em] (n1_1) at (0, -2) {*create\\CoroutineImpl*};
    \node[note, text width=10em] (n1_2) at (0, -4.5) {\texttt{CoroutineImpl::\\resumeWith(Unit)}};
    \node[whiteNote, text width=9em] (n1_3) at (0, -9.5) {return\\\texttt{COROUTINE\_SUSPENDED}};
    \node[whiteNote] (n1_4) at (0, -14.5) {return\\result};

% -------------------------
% Column 2 Nodes
% -------------------------
    \node[pinkNote] (n2_1) at (6, -4.5) {\texttt{doResume()}};
    \node[note] (n2_2) at (6, -9.5) {coroutine suspended};
    \node[note] (n2_3) at (6, -14.5) {resume\\completion\\with result};

% -------------------------
% Column 3 Nodes
% -------------------------
    \node[note] (n3_1) at (12, -4.5) {*execute\\suspend\\fun code*};
    \node[pinkNote, minimum height=2em, text width=10em] (n3_2) at (12, -7) {\texttt{suspendCoroutine-\\UninterceptedOr-\\Return(suspendLambda)}};
    \node[note] (n3_3) at (12, -12) {*continue\\execution*};
    \node[note] (n3_4) at (12, -14.5) {return\\result};

% -------------------------
% Column 4 Nodes
% -------------------------
    \node[whiteNote] (n4_1) at (18, -7) {*execute\\suspend\\lambda*};

    % The branching paths below the lambda execution
    \node[whiteNote, text width=6.5em] (n4_2) at (16, -9.5) {result is\\\texttt{COROUTINE\_\\SUSPENDED}};
    \node[whiteNote, text width=6.5em] (n4_3) at (20, -9.5) {result is\\a value};

% -------------------------
% Background Column Boxes
% -------------------------
    \begin{pgfonlayer}{background}
        % Using dummy coordinates to ensure uniform top/bottom padding
        \coordinate (c1_top) at (0, 0.8);
        \coordinate (c1_bot) at (0, -15.5);
        \node[bgBox, fit=(t1) (n1_1) (n1_4) (c1_top) (c1_bot)] {};

        \coordinate (c4_top) at (18, 0.8);
        \coordinate (c4_bot) at (18, -15.5);
        \node[bgBox, fit=(t4) (n4_1) (n4_2) (n4_3) (c4_top) (c4_bot)] {};
    \end{pgfonlayer}

% -------------------------
% Arrows & Routing
% -------------------------
    % Downward Start
    \draw[arrow] (n1_1) -- (n1_2);

    % Rightward step-in
    \draw[arrow] (n1_2) -- (n2_1);
    \draw[arrow] (n2_1) -- (n3_1);

    % Down to suspendCoroutine call
    \draw[arrow] (n3_1) -- (n3_2);
    \draw[arrow] (n3_2) -- (n4_1);

    % Curving Fork from lambda execution
    \draw[arrow] (n4_1.south) to[out=270, in=90] (n4_2.north);
    \draw[arrow] (n4_1.south) to[out=270, in=90] (n4_3.north);

    % Leftward flow (Suspension path)
    \draw[arrow] (n4_2) -- (n2_2);
    \draw[arrow] (n2_2) -- (n1_3);

    % Long curved flow (Value path)
    \draw[arrow] (n4_3.south) to[out=270, in=0] (n3_3.east);

    % Downward return to caller
    \draw[arrow] (n3_3) -- (n3_4);

    % Leftward flow (Resume path)
    \draw[arrow] (n3_4) -- (n2_3);
    \draw[arrow] (n2_3) -- (n1_4);
\end{tikzpicture}
